% Kurzbeschreibung mail_candidate.py
% Maximal 30 Zeilen, Stichpunkte / kurze Sätze

\textbf{mail\_candidate.py -- E-Mail-Generierung für Kandidaten}

Generiert personalisierte Recruiting-E-Mails für Kandidaten auf Basis von Matching-Ergebnissen.
Nutzt die im Ordner \texttt{results} gespeicherten Excel-Matching-Dateien als Eingangsdaten.
Lädt aus der Excel-Datei Kandidaten- und Job-Informationen (inkl. Gehalt und Stellenbeschreibung).
Prüft, ob der lokale Ollama-Server läuft und welche Modelle verfügbar sind.
Bietet verschiedene LLM-Modelle an (z.B. \texttt{mistral:latest}, \texttt{llama3.2:3b}, \texttt{phi3:3.8b}).
Erlaubt alternativ einen reinen Template-Modus ohne LLM.
Baut für jeden Kandidaten einen Prompt, der Name, Qualifikation, Position, Ort, Gehalt und relevante Skills enthält.
Ruft die Ollama-API auf und erwartet eine Antwort mit \glqq BETREFF:\grqq{} und \glqq MAIL:\grqq{}-Abschnitt.
Parsed aus der LLM-Antwort Betreffzeile und E-Mail-Text; bei Abweichungen wird ein Fallback-Betreff verwendet.
Enthält eine Fallback-Funktion, die ohne LLM eine sinnvolle Standard-E-Mail mit Gehaltsangaben erzeugt.
Hebt in der E-Mail übereinstimmende Skills zwischen Stelle und Kandidat hervor, wenn diese bekannt sind.
Speichert alle erzeugten E-Mails in einzelnen Textdateien im Verzeichnis \texttt{results/job\_ID/}.
Jede Datei enthält Kopfzeilen mit Job-ID, Kandidaten-ID, Name, E-Mail-Adresse und Score.
Gibt nach Abschluss ein Beispiel der ersten E-Mail auf der Konsole aus.
Kann interaktiv im Terminal ausgeführt werden (Auswahl der Matching-Datei und des Modells).

